\section{Conclusion}
\label{sec:9}

This paper presented Software-as-a-Graph (SaG), a pre-deployment framework that derives explicit dependency graphs from publish--subscribe deployment manifests and supports analytical, hybrid, and learned approaches to cascade-impact ranking. Using twelve synthetic architectures, three failure simulators and five stylized models of open-source systems, the study examined when graph learning improves on rankings derived directly from explicit architectural dependencies. The primary registered contrast was null, and the findings below are registered secondary or exploratory.

For the reachability and queue-flow simulators, explicit dependency representations provided most of the predictive signal. Attention-based learners benefited from the derived dependency graph, and that gain survives a control for edge direction,
and the removal of every feature that computes part of the reachability simulator,
but simple analytical rankings built on the same dependencies were not distinguishable from, or exceeded, every learned model: counting direct dependents on the reachability simulator, and a rate-weighted first-order expansion, computed in milliseconds, on the queue-flow simulator, whose labels took $12.7$ CPU-hours to generate.
Learners that start from the rate-weighted approximation matched it at best.

Learning added value only relative to a weak comparator: hybrids beat the training-free baseline, which they nest, but not their own base learners, and learned models transferred zero-shot better than that baseline, while dependency-based references remained stronger and ordering among components that actually propagate failures remained weak. For researchers, the results show why graph-learning approaches should be evaluated against analytical alternatives aligned with the simulator, and why prediction should be distinguished from restatement of a simulator's rule. For practitioners, training-free dependency analysis is effective and inexpensive support for prioritizing architectural review. Learned models can be expected to add value only where no aligned analytical approximation exists; all three simulators here admit one, so this corpus could not exercise that condition.

These findings concern simulated failures on synthetic architectures and single-modeler system models, and their validation against real outages is the natural next step (Section~\ref{sec:limitations}). More broadly, the study suggests that the central challenge in architecture-level cascade-impact analysis is not selecting increasingly sophisticated learning algorithms, but deriving meaningful dependency representations and understanding when learning adds value beyond them. By making this comparison explicit and by releasing a reproducible benchmark with reconciled results, Software-as-a-Graph provides a foundation for more principled evaluation of analytical and learned methods for pre-deployment architecture analysis.
